\documentclass[a4paper,11pt]{article}

\usepackage[a4paper,margin=2cm]{geometry}
\usepackage[T1]{fontenc}
\usepackage[utf8]{inputenc}
\usepackage[ngerman,english]{babel}
\usepackage{lmodern}
\usepackage{microtype}
\usepackage{xcolor}
\usepackage{parskip}
\usepackage{enumitem}
\usepackage[most]{tcolorbox}
\usepackage{titlesec}
\usepackage[hidelinks]{hyperref}

\definecolor{HeaderBlueDark}{RGB}{70,110,170}
\definecolor{RowBlueLight}{RGB}{235,243,255}
\definecolor{RowBlueDark}{RGB}{220,233,250}
\definecolor{SoftGray}{RGB}{90,90,90}

\setlength{\parindent}{0pt}
\setlist[itemize]{leftmargin=1.4em,itemsep=0.35em,topsep=0.35em}
\linespread{1.06}

\titleformat{\section}[block]
{\Large\bfseries\color{HeaderBlueDark}}
{}{0pt}{}
\titlespacing{\section}{0pt}{1.3em}{0.8em}

\titleformat{\subsection}[block]
{\large\bfseries\color{HeaderBlueDark}}
{}{0pt}{}
\titlespacing{\subsection}{0pt}{1.0em}{0.6em}

\newtcolorbox{exbox}{enhanced,breakable,colback=RowBlueLight,colframe=RowBlueDark,boxrule=0.4pt,arc=1.5mm,left=1.2mm,right=1.2mm,top=1mm,bottom=1mm,title={Beispiele \textnormal{(Examples)}},fonttitle=\bfseries,colbacktitle=RowBlueDark,coltitle=black}

\NewDocumentEnvironment{grammarentry}{mm+b}{%
    \begin{tcolorbox}[enhanced,breakable,colback=white,colframe=HeaderBlueDark,boxrule=0.6pt,arc=1.5mm,left=1.5mm,right=1.5mm,top=1mm,bottom=1mm,colbacktitle=HeaderBlueDark,coltitle=white,fonttitle=\bfseries,title={#1\hfill\normalfont\footnotesize #2}]
    #3
    \end{tcolorbox}
}{}

\newcommand{\GermanText}[1]{\textbf{Deutsch: }#1\par}
\newcommand{\EnglishText}[1]{{\small\color{SoftGray}\textbf{English: }#1\par}}
\newcommand{\ExampleItem}[2]{\item \textbf{#1}\\{\small\color{SoftGray}#2}}

\title{Rückmeldung geben}
\date{}

\begin{document}
    \maketitle
    \tableofcontents
    \newpage
    
    
    \section{Bedeutung}
        
        \begin{grammarentry}{Rückmeldung geben}{to give feedback}
            \GermanText{Rückmeldung geben bedeutet: jemandem sagen, wie etwas war, was gut war und was man verbessern kann.}
            \EnglishText{To give feedback means to tell someone how something was, what was good, and what can be improved.}
            
            \begin{exbox}
                \begin{itemize}
                    \ExampleItem{Ich möchte dir eine kurze Rückmeldung geben.}{I would like to give you brief feedback.}
                    \ExampleItem{Danke für deine Rückmeldung.}{Thank you for your feedback.}
                    \ExampleItem{Kannst du mir bitte Rückmeldung zu meinem Text geben?}{Can you please give me feedback on my text?}
                \end{itemize}
            \end{exbox}
        \end{grammarentry}
    
    
    \section{Grundprinzip}
        
        \begin{grammarentry}{Gute Rückmeldung}{good feedback}
            \GermanText{Gute Rückmeldung ist klar, konkret und hilfreich. Sie soll nicht nur sagen, dass etwas falsch ist, sondern auch zeigen, wie man es besser machen kann.}
            \EnglishText{Good feedback is clear, specific, and helpful. It should not only say that something is wrong, but also show how it can be improved.}
            
            \begin{itemize}
                \item zuerst verstehen, was die Person sagen wollte
                \item dann den Fehler oder das Problem zeigen
                \item danach eine bessere Formulierung anbieten
                \item am Ende kurz erklären, warum die neue Form besser ist
            \end{itemize}
        \end{grammarentry}
    
    
    \section{Direkte Rückmeldung}
        
        \begin{grammarentry}{Die direkte Methode}{the direct method}
            \GermanText{Bei der direkten Methode sagt man klar, was falsch ist, und gibt sofort die korrekte Form.}
            \EnglishText{With the direct method, you clearly say what is wrong and immediately give the correct form.}
            
            \begin{exbox}
                \begin{itemize}
                    \ExampleItem{„nach die Bauernhof“ ist falsch. Richtig ist: „auf den Bauernhof“.}{“nach die Bauernhof” is wrong. Correct is: “auf den Bauernhof”.}
                    \ExampleItem{Der Artikel ist hier falsch. Man sagt: der Bauernhof.}{The article is wrong here. You say: der Bauernhof.}
                    \ExampleItem{Die Präposition passt hier nicht. Besser ist: auf den Bauernhof gehen.}{The preposition does not fit here. Better is: auf den Bauernhof gehen.}
                \end{itemize}
            \end{exbox}
        \end{grammarentry}
    
    
    \section{Freundliche Rückmeldung}
        
        \begin{grammarentry}{Die freundliche Methode}{the friendly method}
            \GermanText{Bei der freundlichen Methode beginnt man mit etwas Positivem und korrigiert danach ruhig und respektvoll.}
            \EnglishText{With the friendly method, you begin with something positive and then correct calmly and respectfully.}
            
            \begin{exbox}
                \begin{itemize}
                    \ExampleItem{Ich verstehe, was du meinst. Natürlicher sagt man aber: „Ich gehe auf den Bauernhof.“}{I understand what you mean. But more naturally, you say: “Ich gehe auf den Bauernhof.”}
                    \ExampleItem{Deine Idee ist klar, aber die Grammatik braucht noch eine kleine Korrektur.}{Your idea is clear, but the grammar still needs a small correction.}
                    \ExampleItem{Du hast es schön beschrieben. Nur die Präposition würde ich ändern.}{You described it well. I would only change the preposition.}
                \end{itemize}
            \end{exbox}
        \end{grammarentry}
    
    
    \section{Kritische Rückmeldung}
        
        \begin{grammarentry}{Die kritische Methode}{the critical method}
            \GermanText{Kritische Rückmeldung zeigt deutlich, wo ein Problem ist. Sie sollte trotzdem sachlich und nicht beleidigend sein.}
            \EnglishText{Critical feedback clearly shows where there is a problem. However, it should still be factual and not insulting.}
            
            \begin{exbox}
                \begin{itemize}
                    \ExampleItem{Dieser Satz ist grammatisch nicht korrekt.}{This sentence is grammatically not correct.}
                    \ExampleItem{Die Wortstellung ist hier falsch.}{The word order is wrong here.}
                    \ExampleItem{Dieser Ausdruck klingt im Deutschen nicht natürlich.}{This expression does not sound natural in German.}
                    \ExampleItem{Der Satz ist verständlich, aber er braucht eine bessere Struktur.}{The sentence is understandable, but it needs a better structure.}
                \end{itemize}
            \end{exbox}
        \end{grammarentry}
    
    
    \section{Die Sandwich-Methode}
        
        \begin{grammarentry}{Lob, Korrektur, Ermutigung}{praise, correction, encouragement}
            \GermanText{Bei der Sandwich-Methode gibt man zuerst Lob, dann eine Korrektur und am Ende eine Ermutigung.}
            \EnglishText{With the sandwich method, you first give praise, then a correction, and at the end encouragement.}
            
            \begin{exbox}
                \begin{itemize}
                    \ExampleItem{Deine Idee ist gut und ich verstehe deinen Satz.}{Your idea is good and I understand your sentence.}
                    \ExampleItem{Besser wäre aber: „Ich gehe auf den Bauernhof.“}{But better would be: “Ich gehe auf den Bauernhof.”}
                    \ExampleItem{Wenn du auf die Präpositionen achtest, klingt dein Deutsch viel natürlicher.}{If you pay attention to the prepositions, your German will sound much more natural.}
                \end{itemize}
            \end{exbox}
        \end{grammarentry}
    
    
    \section{Die fragende Methode}
        
        \begin{grammarentry}{Korrektur durch Fragen}{correction through questions}
            \GermanText{Bei dieser Methode korrigiert man nicht sofort direkt, sondern stellt zuerst eine Frage. Das ist besonders freundlich und hilft der anderen Person, selbst nachzudenken.}
            \EnglishText{With this method, you do not correct directly immediately, but first ask a question. This is especially friendly and helps the other person think for themselves.}
            
            \begin{exbox}
                \begin{itemize}
                    \ExampleItem{Meinst du, dass du auf dem Bauernhof arbeitest oder dass du zum Bauernhof gehst?}{Do you mean that you work on the farm or that you go to the farm?}
                    \ExampleItem{Willst du hier eine Bewegung oder eine Position ausdrücken?}{Do you want to express movement or position here?}
                    \ExampleItem{Meinst du „unter der Sonne“ oder „in der Sonne“?}{Do you mean “under the sun” or “in the sun”?}
                \end{itemize}
            \end{exbox}
        \end{grammarentry}
    
    
    \section{Die erklärende Methode}
        
        \begin{grammarentry}{Regel plus Beispiel}{rule plus example}
            \GermanText{Bei der erklärenden Methode gibt man zuerst die Regel und danach ein oder zwei Beispiele.}
            \EnglishText{With the explanatory method, you first give the rule and then one or two examples.}
            
            \begin{exbox}
                \begin{itemize}
                    \ExampleItem{Bei Bewegung benutzt man den Akkusativ: Ich gehe auf den Bauernhof.}{For movement, you use the accusative: I go to the farm.}
                    \ExampleItem{Bei Position benutzt man den Dativ: Ich arbeite auf dem Bauernhof.}{For position, you use the dative: I work on the farm.}
                    \ExampleItem{Nach „um“ plus Infinitiv braucht man „zu“: Ich gehe, um zu arbeiten.}{After “um” plus infinitive, you need “zu”: I go in order to work.}
                \end{itemize}
            \end{exbox}
        \end{grammarentry}
    
    
    \section{Kurze Redemittel}
        
        \begin{grammarentry}{Nützliche Sätze}{useful sentences}
            \GermanText{Diese Sätze kann man benutzen, wenn man schnell und klar Rückmeldung geben möchte.}
            \EnglishText{These sentences can be used when you want to give quick and clear feedback.}
            
            \begin{exbox}
                \begin{itemize}
                    \ExampleItem{Das ist richtig.}{That is correct.}
                    \ExampleItem{Das ist nicht ganz richtig.}{That is not completely correct.}
                    \ExampleItem{Das klingt natürlich.}{That sounds natural.}
                    \ExampleItem{Das klingt nicht natürlich.}{That does not sound natural.}
                    \ExampleItem{Besser sagt man: ...}{It is better to say: ...}
                    \ExampleItem{Natürlicher klingt: ...}{More natural is: ...}
                    \ExampleItem{Im Alltag sagt man eher: ...}{In everyday language, people usually say: ...}
                    \ExampleItem{Grammatisch ist das möglich, aber es klingt ungewöhnlich.}{Grammatically, that is possible, but it sounds unusual.}
                \end{itemize}
            \end{exbox}
        \end{grammarentry}
    
    
    \section{Rückmeldung zu Texten}
        
        \begin{grammarentry}{Textkorrektur}{text correction}
            \GermanText{Wenn man einen Text korrigiert, sollte man zuerst die korrigierte Version zeigen. Danach kann man die Fehler Satz für Satz erklären.}
            \EnglishText{When correcting a text, one should first show the corrected version. After that, the mistakes can be explained sentence by sentence.}
            
            \begin{exbox}
                \begin{itemize}
                    \ExampleItem{Korrigierter Satz: Ich gehe auf den Bauernhof, um dort zu arbeiten.}{Corrected sentence: I go to the farm in order to work there.}
                    \ExampleItem{Fehler: „nach die Bauernhof“ ist falsch.}{Mistake: “nach die Bauernhof” is wrong.}
                    \ExampleItem{Erklärung: Man sagt „auf den Bauernhof“, weil der Bauernhof als Gelände verstanden wird.}{Explanation: You say “auf den Bauernhof” because the farm is understood as an area.}
                \end{itemize}
            \end{exbox}
        \end{grammarentry}
    
    
    \section{Sehr natürliche Rückmeldung}
        
        \begin{grammarentry}{Alltagssprache}{everyday language}
            \GermanText{Diese Sätze klingen im Alltag natürlich und freundlich.}
            \EnglishText{These sentences sound natural and friendly in everyday language.}
            
            \begin{exbox}
                \begin{itemize}
                    \ExampleItem{Ja, ich verstehe dich. Ich würde es nur ein bisschen anders sagen.}{Yes, I understand you. I would just say it a little differently.}
                    \ExampleItem{Der Satz ist fast richtig. Nur ein kleines Detail fehlt.}{The sentence is almost correct. Only one small detail is missing.}
                    \ExampleItem{Das kannst du sagen, aber diese Version klingt natürlicher.}{You can say that, but this version sounds more natural.}
                    \ExampleItem{Sehr gut, genau so sagt man das.}{Very good, that is exactly how you say it.}
                    \ExampleItem{Das war schon viel besser als vorher.}{That was already much better than before.}
                \end{itemize}
            \end{exbox}
        \end{grammarentry}

\end{document}